\documentclass[a4paper, 11pt]{article}

\newcommand{\projectname}{Digital håndtering af færgebilletter}
\newcommand{\student}{Marcus Haukelid Larsen}
\newcommand{\supervisor}{Lars Thise Pedersen og Simon Hoxer Bønding}

% --- Language and Encoding ------------------------------------------------------------------------
\usepackage[danish]{babel}

% --- Layout ---------------------------------------------------------------------------------------
\usepackage{parskip}
\usepackage{float}
\usepackage{graphicx}
\usepackage{svg}

% --- Code -----------------------------------------------------------------------------------------
\usepackage{listings}
\lstset{
  basicstyle=\ttfamily\small,
  showspaces=false,
  showstringspaces=false,
  breaklines=true,
  frame=single
}

% --- Bibliography ---------------------------------------------------------------------------------
\usepackage[
  backend=biber,
  style=numeric
]{biblatex}

% --- Links ----------------------------------------------------------------------------------------
\usepackage[hidelinks]{hyperref}

\addbibresource{process-report.bib}

\newcommand{\projecttype}{Procesrapport}

\begin{document}

\pdfbookmark[1]{Forside}{frontpage}

\begin{center}
  {\Huge \textbf{\projectname}}

  \vspace*{\fill}

  {\Large \textbf{\student}}

  \vspace*{\fill}

  {\large \textbf{\projecttype}}
\end{center}

\thispagestyle{empty}

\newpage
\begin{center}
  {\large \textbf{Elev:} \student}
  \vspace*{\fill}

  {\large \textbf{Projektnavn:} \projectname}
  \vspace*{\fill}

  {\large \textbf{Uddannelsessted:} Techcollege, Struervej 70, 9220 Aalborg Øst, Denmark}
  \vspace*{\fill}

  {\large \textbf{Elevplads:} RTX A/S, Strømmen 6, 9400 Nørresundby, Denmark}
  \vspace*{\fill}

  {\large \textbf{Projektperiode:} 31/08-2026 - 29/09-2026}
  \vspace*{\fill}

  {\large \textbf{Afleveringsdato:} 24/09-2026}
  \vspace*{\fill}

  {\large \textbf{Fremlæggelsesdato:} 30/09-2026}
  \vspace*{\fill}

  {\large \textbf{Vejledere:} \supervisors}
  \vspace*{\fill}

  {\large \textbf{Underskrift:}}
  \vspace*{\fill}

  \begin{tabular}{p{6cm} p{6cm}}
    \hrulefill & \hrulefill \\
    Vejledere (\supervisors) & Elev (\student)
  \end{tabular}
\end{center}

\thispagestyle{empty}


\newpage
\tableofcontents

\newpage

\section{Læsevejledning}

\section{Indledning}

\section{Casebeskrivelse}

\section{Problemformulering}

\section{Projektplanlægning}

\subsection{Tidsplan}

\section{Sprogvalg}

Projektets analyse og dokumentation blev udarbejdet på dansk, mens kildekode, filnavne og øvrige
tekniske elementer såsom API-specifikationen blev skrevet på engelsk. Dette sikrede et ensartet
sprog i kildekoden og den tekniske kontekst, da engelsk er det primære sprog inden for
programmering.

\section{Metode- og teknologivalg}

\section{Udviklingsproces}

\subsection{Analyse og design}

\subsubsection{Domænemodel}

\begin{figure}[H]
  \centering
  \includegraphics[width=1\textwidth]{domain-model.png}
  \caption{Domænemodel}
  \label{fig:domain-model}
\end{figure}

\subsubsection{Datamodel}

\begin{figure}[H]
  \centering
  \includegraphics[width=1\textwidth]{database-model.png}
  \caption{Database model}
  \label{fig:database-model}
\end{figure}

\subsubsection{Brugerflow}

\begin{figure}[H]
  \centering
  \includegraphics[width=1\textwidth]{reservation-flow.png}
  \caption{Reservation flow}
  \label{fig:reservation-flow}
\end{figure}

\begin{figure}[H]
  \centering
  \includegraphics[width=1\textwidth]{validation-flow.png}
  \caption{Validerings flow}
  \label{fig:validerings-flow}
\end{figure}

\subsection{Database og API}

\subsection{Server}

\subsection{Klient}

\subsection{Test og kvalitetssikring}

\subsection{Konfigurationsstyring}

\subsection{CI/CD og hosting}

\section{Afgrænsning}

\section{Realiseret tidsplan}

\section{Evaluering}

\subsection{Afvigelser fra projektplanen}

\subsection{Erfaringer}

\section{Konklusion}

\printbibliography[title={Referencer}]

\appendix
\section{Logbog}

\subsection{Dag 1 -- 31/08-2026}

Logbogen blev oprettet, og det første udkast til designdokumentet blev udarbejdet.

\subsection{Dag 2 -- 01/09-2026}

Specifikationen og modellerne blev finpudset på baggrund af feedback om manglende information.

\subsection{Dag 3 -- 02/09-2026}

Specifikationen og modellerne blev finpudset på baggrund af feedback om manglende information.
API-beskrivelsen blev opdateret, så den stemmer overens med modellerne, og der blev udarbejdet
brugerflows.

\subsection{Dag 4 -- 03/09-2026}

Der blev udarbejdet en skitse af brugergrænsefladen.

\subsection{Dag 5 -- 04/09-2026}

Udviklingen af serveren blev påbegyndt. Undervejs blev der identificeret en designfejl i
API-specifikationen, som var blevet udarbejdet på et for lavt abstraktionsniveau. Ideelt set burde
systemet have både et high-level og et low-level API, men på grund af den begrænsede tid blev der
udelukkende fokuseret på high-level API-designet.

\subsection{Dag 6 -- 05/09-2026}

API'et blev finpudset, og de resterende slutpunkter blev implementeret. Herefter blev der testet
manuelt for at kontrollere, at oprettelse, læsning og opdatering fungerer som forventet.

\subsection{Dag 7 -- 07/09-2026}

Udviklingen af klienten blev påbegyndt, og projektstrukturen blev tilpasset ved at omdøbe public
til asset. README-filen blev desuden udvidet, og henvisningerne til diagrammer blev opdateret.

\subsection{Dag 8 -- 08/09-2026}

Navigationen i rejsevælgeren blev implementeret, hvorefter dens placering og udseende blev
finpudset. Udrejsepanelet blev udarbejdet med data fra API'et, som samtidig blev tilpasset og
finpudset.

\subsection{Dag 9 -- 09/09-2026}

API'et blev udvidet med yderligere data, og klientkoden blev omstruktureret for at gøre den mere
overskuelig.

\subsection{Dag 10 -- 10/09-2026}

Casebeskrivelsen og problemformuleringen blev omskrevet på baggrund af feedback fra vejleder, så
det fremgår tydeligt, at systemet omhandler færgebilletter. Ydermere blev billetvælgeren
udarbejdet, så kunden kan vælge og oprette de ønskede billetter.

\subsection{Dag 11 -- 11/09-2026}

Klientkoden blev finpudset og struktureret yderligere. Der blev desuden arbejdet videre med
billetvælgeren, herunder håndtering af dropdown-menuer og valg af billet- og køretøjsvarianter.

\subsection{Dag 12 -- 14/09-2026}

Datahåndteringen i billetvælgeren blev videreudviklet og finpudset. Billetterne blev
omstruktureret, så den samme datastruktur benyttes ved både hentning og oprettelse af billetter.
Derudover blev indlæsningen af data fra API'et udvidet og tilpasset den nye datastruktur.

\subsection{Dag 13 -- 15/09-2026}

Der blev udarbejdet en kalender-widget til visning og valg af afgange. Afgangene blev grupperet
efter dato og placeret i kalenderen ud fra uge og ugedag. Derudover blev API'et justeret, så de
nødvendige data for en afgang er korrekte.

\subsection{Dag 14 -- 16/09-2026}

Afgangsvælgeren blev udarbejdet, så afgange kan vises og vælges for den valgte dag. Derudover
blev klienten udvidet med mulighed for at oprette billetter gennem API'et.

\subsection{Dag 15 -- 17/09-2026}

Produktet blev klargjort til hosting, og der blev foretaget den nødvendige finpudsning og
fejlfinding i forbindelse med deployment. Derudover blev der udarbejdet unit tests til serveren
for at verificere API'ets funktionalitet og databasehåndtering.

\subsection{Dag 16 -- 18/09-2026}

Proces- og produktrapporten blev påbegyndt. Bygge scriptet blev udvidet til at generere rapporterne,
og CI-konfigurationen blev udarbejdet med de nødvendige LaTeX-afhængigheder, så rapporterne
automatisk kan bygges ved en Github release.


\end{document}